\documentclass[11pt,a4paper]{article}
\usepackage[UTF8]{ctex}
\usepackage{amsmath,amssymb,bm}
\usepackage{booktabs,longtable,array}
\usepackage{enumitem,fancyhdr,geometry,needspace}
\usepackage{xcolor,listings}
\usepackage[most]{tcolorbox}
\usepackage{xurl,hyperref}
\geometry{left=2.15cm,right=2.15cm,top=2.05cm,bottom=2.05cm}
\setlength{\headheight}{14pt}
\setlength{\parindent}{2em}
\setlength{\parskip}{0.28em}
\setlength{\emergencystretch}{2em}
\setlength{\LTleft}{0pt}
\setlength{\LTright}{\fill}
\setlist[itemize]{leftmargin=2em,itemsep=0.25em,topsep=0.3em}
\setlist[enumerate]{leftmargin=2.3em,itemsep=0.3em,topsep=0.3em}
\definecolor{courseblue}{HTML}{2457A6}
\definecolor{coursegreen}{HTML}{18794E}
\definecolor{codered}{HTML}{A13D2D}
\definecolor{codegray}{HTML}{F4F6F8}
\hypersetup{colorlinks=true,linkcolor=courseblue,urlcolor=courseblue,pdftitle={LAB1 线性回归与正则化},pdfauthor={中国科学技术大学}}
\newcommand{\code}[1]{{\small\nolinkurl{#1}}}
\lstdefinestyle{coursecode}{basicstyle=\ttfamily\small,backgroundcolor=\color{codegray},frame=single,rulecolor=\color{black!12},breaklines=true,columns=fullflexible,keepspaces=true,showstringspaces=false,aboveskip=0.55em,belowskip=0.55em}
\lstset{style=coursecode}
\newtcolorbox{importantbox}{enhanced,breakable,colback=courseblue!4,colframe=courseblue!65,boxrule=0.7pt,arc=1.5mm,left=1.5mm,right=1.5mm,top=1.2mm,bottom=1.2mm}
\newtcolorbox{tipbox}{enhanced,breakable,colback=coursegreen!4,colframe=coursegreen!65,boxrule=0.7pt,arc=1.5mm,left=1.5mm,right=1.5mm,top=1.2mm,bottom=1.2mm}
\pagestyle{fancy}
\fancyhf{}
\lhead{AI3002 人工智能与机器学习基础}
\rhead{LAB1：线性回归与正则化}
\cfoot{\thepage}
\renewcommand{\headrulewidth}{0.4pt}
\fancypagestyle{plain}{\fancyhf{}\lhead{AI3002 人工智能与机器学习基础}\rhead{LAB1：线性回归与正则化}\cfoot{\thepage}\renewcommand{\headrulewidth}{0.4pt}}
\title{\vspace{-0.8cm}\textbf{人工智能与机器学习基础\\LAB1：线性回归与正则化}}
\author{中国科学技术大学}
\date{截止时间：2026 年 10 月 19 日 06:00（北京时间）}
\begin{document}
\maketitle
\thispagestyle{fancy}
\setcounter{section}{-1}

\Needspace{7\baselineskip}\section*{环境创建（不计分）}
沿用 LAB0 已安装的 Miniforge 和 VS Code，为本实验单独创建名为 \code{ai3002-lab1} 的 Conda 环境，不在 LAB0 的 \code{ai3002} 环境中安装或运行。使用 Python 3.11，CPU 即可，无需安装 PyTorch 或配置 GPU。本部分不计分，不提交 OJ 验证码。


\Needspace{5\baselineskip}\subsection*{Windows}
打开 PowerShell，进入你保存的 LAB1 目录。若课程仓库位于 LAB0 示例位置，可执行下面的路径命令；实际位置不同则替换第一行。


\par\addvspace{0.65\baselineskip}\noindent\begin{minipage}{\linewidth}
\begin{lstlisting}[aboveskip=0pt,belowskip=0pt]
cd $HOME\ALML_public\LABs\LAB1
conda env create -f environment.yml
conda activate ai3002-lab1
python --version
where.exe python
python -c "import sys; print(sys.executable)"
\end{lstlisting}
\end{minipage}\par\addvspace{0.65\baselineskip}
若 PowerShell 无法识别 \code{conda}，按 LAB0 的方法打开 Miniforge Prompt，执行 {\small\texttt{conda init powershell}}，关闭并重新打开 PowerShell 后再创建环境。


\Needspace{5\baselineskip}\subsection*{macOS}
打开 Terminal，进入实际的 LAB1 目录；下面路径沿用 LAB0 示例。


\par\addvspace{0.65\baselineskip}\noindent\begin{minipage}{\linewidth}
\begin{lstlisting}[aboveskip=0pt,belowskip=0pt]
cd "$HOME/ALML_public/LABs/LAB1"
conda env create -f environment.yml
conda activate ai3002-lab1
python --version
which python
python -c "import sys; print(sys.executable)"
\end{lstlisting}
\end{minipage}\par\addvspace{0.65\baselineskip}
若无法识别 \code{conda}，参照 LAB0 初始化当前 Shell，并重新打开终端。不要把 Windows 的路径命令复制到 macOS 终端。


\Needspace{5\baselineskip}\subsection*{选择解释器并检查}
在 VS Code 中打开 LAB1 文件夹，Windows 按 \code{Ctrl+Shift+P}、macOS 按 \code{Cmd+Shift+P}，执行 {\small\texttt{Python: Select Interpreter}}，选择 \code{ai3002-lab1} 对应的 Python 3.11。新建集成终端，确认环境已激活，再在 LAB1 根目录执行：


\par\addvspace{0.65\baselineskip}\noindent\begin{minipage}{\linewidth}
\begin{lstlisting}[aboveskip=0pt,belowskip=0pt]
python src/verify_env.py
\end{lstlisting}
\end{minipage}\par\addvspace{0.65\baselineskip}
脚本检查新环境名称、Python 版本、NumPy、scikit-learn 和 Matplotlib，输出 {\small\texttt{LAB1 environment ready}}，并在 \code{results/} 下生成环境记录和绘图检查文件。环境记录仅用于排错，不占分。若解释器仍指向 \code{base} 或 \code{ai3002}，重新选择解释器并激活新环境，不要修改检查脚本。

后续每次实验先执行 {\small\texttt{conda activate ai3002-lab1}}，所有命令均在 LAB1 根目录运行，无需进入 \code{src/}。用 {\small\texttt{conda info --envs}} 查看环境；已创建成功的环境不必重复创建，用完执行 {\small\texttt{conda deactivate}} 即可。


\par\addvspace{0.65\baselineskip}\noindent\begin{minipage}{\linewidth}
\begin{lstlisting}[aboveskip=0pt,belowskip=0pt]
LAB1/
|-- lab1.pdf
|-- README.md
|-- environment.yml
|-- requirements.txt
|-- src/
|   |-- submission.py
|   |-- run.py
|   |-- supplied.py
|   |-- config_protocol.py
|   |-- benchmark.py
|   `-- verify_env.py
|-- data/
|   `-- dev.csv
`-- results/
\end{lstlisting}
\end{minipage}\par\addvspace{0.65\baselineskip}

\Needspace{7\baselineskip}\section{实验介绍}
本实验用线性模型预测设备一次运行的能耗，熟悉数据预处理、模型实现、训练、验证和测试的完整流程。预测发生在运行开始之前，模型将用于训练数据中没有出现过的新设备。


\begin{importantbox}
个人完成，CPU 即可，预计 8--10 小时，满分 100 分：OJ 系统评分 50 分（代码 40 分、最终性能 10 分），实验报告评分 50 分。必做内容为第 1、2 部分和第 4 部分的四个问题。环境与提交要求见 \code{README.md}。补全 \code{src/submission.py} 中的七个计算函数，调参后填写 \code{get_config}()，并完成学习率与正则化参数的两轮调整；数据读取、固定划分、实验循环、画图和模型保存由 \code{src/run.py} 提供。

\end{importantbox}

\Needspace{5\baselineskip}\subsection*{0.1 数据与文件}
\code{data/dev.csv} 包含 60 台设备的 600 条记录，每台设备 10 条。\code{supplied.load_development()} 返回按 \code{row_id} 排序的数组。原始特征 \code{X} 的形状为 $(600,29)$，标签 \code{y} 的形状为 $(600,)$。


{\small\renewcommand{\arraystretch}{1.18}
\begin{longtable}{@{}>{\raggedright\arraybackslash}p{\dimexpr 0.26\linewidth-2\tabcolsep\relax}>{\raggedright\arraybackslash}p{\dimexpr 0.74\linewidth-2\tabcolsep\relax}@{}}
\toprule
\textbf{字段} & \textbf{含义与用途} \\
\midrule
\endfirsthead
\toprule
\textbf{字段} & \textbf{含义与用途} \\
\midrule
\endhead
\midrule
\multicolumn{2}{r}{\footnotesize 接下页}\\
\endfoot
\bottomrule
\endlastfoot
load、temperature、setting & 运行前的负载、温度与设定值；可作特征 \\
\code{load_near}、\code{temperature_milli} & 相关或不同单位的运行前测量；可作特征 \\
\code{nuisance_00} 至 \code{nuisance_23} & 24 个其他运行前测量；可作特征 \\
\code{row_id}、\code{device_id} & 行索引、设备编号；用于匹配和划分，不作特征 \\
\code{post_cycle_meter} & 运行结束后的读数；只在第 2.3 节指定对照中使用 \\
y & 待预测能耗；不能作为输入特征 \\
\end{longtable}}
数据为教学合成数据，无缺失值，不要填补或删除样本。特征顺序由 \code{supplied.SENSOR_NAMES} 指定。本实验使用固定设备划分，不要求交叉验证。


{\small\renewcommand{\arraystretch}{1.18}
\begin{longtable}{@{}>{\raggedright\arraybackslash}p{\dimexpr 0.26\linewidth-2\tabcolsep\relax}>{\raggedright\arraybackslash}p{\dimexpr 0.74\linewidth-2\tabcolsep\relax}@{}}
\toprule
\textbf{文件} & \textbf{作用} \\
\midrule
\endfirsthead
\toprule
\textbf{文件} & \textbf{作用} \\
\midrule
\endhead
\midrule
\multicolumn{2}{r}{\footnotesize 接下页}\\
\endfoot
\bottomrule
\endlastfoot
\code{src/submission.py} & 补全七个待实现函数，是需要修改的代码文件 \\
\code{src/run.py} & 提供检查、比较、最终测试与保存流程，不需要修改 \\
\code{src/supplied.py}、data/ & 提供数据读取和小规模检查样例，不需要修改 \\
results/ & 运行脚本自动生成结果；不要手工编辑结果 \\
\end{longtable}}

\Needspace{5\baselineskip}\subsection*{0.2 固定划分与模型约定}
脚本使用种子 \code{20260918} 打乱排序后的设备编号，前 48 台用于训练，接着 6 台用于验证，最后 6 台用于本地测试；对应 480、60、60 条记录，比例为 $8:1:1$。同一设备不能跨集合。每次拟合的预处理统计量只能从当前训练部分计算。

先在训练集拟合，用验证集选择配置；配置确定后，合并训练集与验证集重训，在本地测试集评价一次；最后用全部 600 条记录重训并在本地保存最终模型。教师另保留新设备数据用于第 6 节的最终性能计分，不要求学生提供该数据的预测成绩。不得根据本地测试成绩继续选择特征或参数。

本实验允许截距，通过特征标准化和标签中心化实现。设当前训练特征为 $F\in\mathbb{R}^{n\times p}$，标签为 $y\in\mathbb{R}^n$。其中 $n$ 为训练样本数，$p$ 为特征数，$i$ 为样本下标，$j$ 为特征下标。定义


\[
\mu_j=\frac1n\sum_{i=1}^n F_{ij},\qquad
s_j=\sqrt{\frac1n\sum_{i=1}^n(F_{ij}-\mu_j)^2},\qquad
\bar y=\frac1n\sum_{i=1}^n y_i.
\]
若 $s_j<10^{-12}$，将该列除数置为 $1$。记标准化矩阵为 $Z$，中心化标签为 $y_c$，系数为 $w\in\mathbb{R}^p$：


\[
Z_{ij}=\frac{F_{ij}-\mu_j}{s_j},\qquad
(y_c)_i=y_i-\bar y,\qquad
\hat y_i=\bar y+\sum_{j=1}^p\frac{F_{ij}-\mu_j}{s_j}w_j.
\]
对验证或测试输入使用已保存的训练均值与标准差，不重新计算；标签只中心化，不除以标准差。截距不受正则化惩罚。以下计算均使用双精度浮点数。


\Needspace{7\baselineskip}\section{实现线性回归（代码 40 分）}
允许使用 NumPy 数组运算、\code{lstsq}、\code{solve} 和 \code{eigvalsh}。这七个函数内不得调用机器学习库的回归器、标准化器、指标函数或自动求导；允许在选做部分另行使用库函数对照。不要求自己实现矩阵分解。


\Needspace{5\baselineskip}\subsection*{1.1 数据预处理（8 分）}


\textbf{（a）实现 {\small\texttt{make\_features(X, basis)}}（3 分）。} \code{linear} 原样保留 29 列；\code{quadratic} 在这 29 列之后，依次追加下列三列，返回 32 列。不要生成其他交叉项，也不要加入常数列。


\[
\mathrm{load}^2,\qquad
\mathrm{load}\,\mathrm{temperature},\qquad
\mathrm{temperature}^2.
\]
\textbf{（b）实现 {\small\texttt{fit\_preprocessor(F, y)}} 和 {\small\texttt{transform(F, prep)}}（5 分）。} 前者返回字典 \code{mu}、\code{scale}、\code{y_mean}，数组形状分别为 $(p,)$、$(p,)$ 和标量；后者只应用已保存的统计量，返回形状与输入相同的矩阵。按第 0.2 节处理常数列，标准差分母使用 $n$。不得修改输入数组。


\Needspace{5\baselineskip}\subsection*{1.2 求解与预测（8 分）}
\textbf{实现 {\small\texttt{solve\_weights(Z, yc, lam)}}。} 参数 \code{lam} 对应正则化强度 $\lambda\geq0$，返回形状为 $(p,)$ 的系数。当前目标为


\[
J_\lambda(w)=\frac1{2n}\|Zw-y_c\|_2^2+
\frac\lambda2\|w\|_2^2.
\]
当 $\lambda=0$，使用 {\small\texttt{np.linalg.lstsq(..., rcond=None)}} 求普通最小二乘解（4 分）。当 $\lambda>0$，利用下式调用 \code{np.linalg.solve}，其中 $I_p$ 为 $p$ 阶单位矩阵（4 分）：


\[
\left(\frac{Z^\top Z}{n}+\lambda I_p\right)w
=\frac{Z^\top y_c}{n}.
\]
预测与截距还原已在 \code{run.py} 中提供。小样例含重复列，普通最小二乘的系数不一定唯一，应检查预测是否一致，不要求自行求逆。


\Needspace{5\baselineskip}\subsection*{1.3 损失与梯度（8 分）}
\textbf{实现 {\small\texttt{loss\_gradient(Z, yc, w, lam)}}。} 返回目标函数标量（3 分）及其关于系数的梯度向量（5 分）。用上一节的目标函数自行写出梯度并实现；报告中列出使用的梯度公式即可，无需重复书面作业的推导。

脚本在固定小样例上以步长 $h=10^{-6}$ 进行中心差分核验，自动记录最大误差。有限差分用于核验，不得作为训练所用的梯度计算方式。


\Needspace{5\baselineskip}\subsection*{1.4 梯度下降训练（8 分）}
\textbf{实现 \code{train_gd}。} 从零系数开始，每次使用全部训练数据，调用自己实现的 \code{loss_gradient}，按下式更新，其中 $t$ 为更新次数，$\eta$ 为传入的学习率：


\[
w^{(t+1)}=w^{(t)}-\eta\nabla J_\lambda(w^{(t)}).
\]
记录初始状态及每次更新后的迭代次数、目标函数值、梯度各分量绝对值的最大值。后者不超过 $10^{-7}$ 时停止，更新次数不得超过传入的 \code{max_iter}；函数默认上限为 20000 次，供第 1 部分精度检查使用，实际实验须服从各阶段传入的预算。为处理过大学习率，若目标或梯度出现非有限值，或目标超过初始目标与 $1$ 中较大者的 $10^6$ 倍，记录该状态并停止。返回最终系数和三列历史数组（更新与停止 5 分，记录 3 分），不得隐藏发散记录。

首先在脚本提供的小样例上用 $\lambda=0.1$ 检查梯度下降实现，检查阶段自动使用以下学习率；$\lambda_{\max}$ 表示对称矩阵的最大特征值。实际设备数据上的学习率调整是第 2.1 节的必做计分任务：


\[
\eta=\frac{1}{\lambda_{\max}(Z^\top Z/n)+0.1}.
\]
执行检查命令，阅读 \code{checks.csv}，查看 \code{loss_curve.png}，核对梯度下降与直接求解的最大预测差不超过 $10^{-5}$。其他检查的容差写在结果文件中。若失败，应定位原因并修正，不要通过修改容差掩盖错误。


\par\addvspace{0.65\baselineskip}\noindent\begin{minipage}{\linewidth}
\begin{lstlisting}[aboveskip=0pt,belowskip=0pt]
python src/run.py --stage check --out results
\end{lstlisting}
\end{minipage}\par\addvspace{0.65\baselineskip}

\Needspace{5\baselineskip}\subsection*{1.5 性能评估（8 分）}
\textbf{实现 {\small\texttt{evaluate(y, pred)}}。} 对形状为 $(m,)$ 的真实值与预测值返回 \code{mse}、\code{rmse}、\code{mae} 三项标量。其中 $m$ 是当前评价样本数，$\hat y_i$ 是预测值：


\[
\mathrm{MSE}=\frac1m\sum_{i=1}^m(\hat y_i-y_i)^2,
\qquad \mathrm{RMSE}=\sqrt{\mathrm{MSE}},
\qquad \mathrm{MAE}=\frac1m\sum_{i=1}^m|\hat y_i-y_i|.
\]
评分为均方误差 3 分、均方根误差 2 分、平均绝对误差 3 分。评价误差不包含正则化项。脚本使用真实值 $(1,2,3,4)$ 和预测值 $(1,2,3,8)$ 检查这三个指标。


\Needspace{7\baselineskip}\section{参数调整与性能比较（实验结果 30 分）}
本部分将学习率调整计 10 分、正则化参数调整计 12 分、评估对照计 4 分、测试保存计 4 分。运行脚本提供实验工具，你需要依据第一轮结果选择新参数，记录“观察到什么、为什么这样调整、调整后如何”，不能只在报告中粘贴默认脚本输出。


\Needspace{5\baselineskip}\subsection*{2.1 调整学习率（10 分）}
固定原设备训练集的 480 条记录、二次特征及 $\lambda=0.1$，本节观察收敛与发散，所有试验从零系数开始，最多更新 2000 次；第 2.2 节起的模型选择与最终训练统一使用 200 次更新预算。令


\[
L=\lambda_{\max}(Z^\top Z/n)+0.1,\qquad \eta=\frac{c}{L}.
\]
其中 $L$ 由当前训练数据计算，$c$ 是学习率倍数。先运行以下命令，比较 $c=0.01$、$1$、$2.1$ 三种情形：


\par\addvspace{0.65\baselineskip}\noindent\begin{minipage}{\linewidth}
\begin{lstlisting}[aboveskip=0pt,belowskip=0pt]
python src/run.py --stage tune-lr --data data --out results
\end{lstlisting}
\end{minipage}\par\addvspace{0.65\baselineskip}
\textbf{（a）初步比较（4 分）。} 阅读 \code{lr_coarse.csv} 和 \code{lr_curves.png}，比较目标函数下降、到达指定精度所需的迭代次数及发散情况。三种试验必须使用相同数据、初始化、正则化参数和迭代预算。第 1.4 节的发散停止规则用于保留失败证据，发散本身不扣分。

\textbf{（b）依据结果再调整一次（4 分）。} 自行选择一个新的倍数 $c_{\mathrm{new}}\in(0,2)$，不能与已试值相同。先在报告中写明选择理由，再运行带 \code{--lr-factor} 的命令，将其与第一轮结果比较。例如下面仅展示命令格式，不指定必须采用的值：


\par\addvspace{0.65\baselineskip}\noindent\begin{minipage}{\linewidth}
\begin{lstlisting}[aboveskip=0pt,belowskip=0pt]
python src/run.py --stage tune-lr --lr-factor 0.5 --out results
\end{lstlisting}
\end{minipage}\par\addvspace{0.65\baselineskip}
\textbf{（c）比较效率并记录（2 分）。} 四组最终记录写入 \code{lr_tuning.csv}，完整历史写入 \code{lr_history.csv}，选择结果写入 \code{lr_config.json}。直接解给出同一训练目标的最优值 $J_*$，脚本比较各次迭代首次满足下式所需的次数：


\[
\left|J_\lambda(w^{(t)})-J_*\right|
\leq 10^{-6}\max(1,|J_*|).
\]
在未发散且达到上述精度的候选中，选择所需次数最少的学习率；次数相同则取较小的倍数。该精度只是统一的比较标准，不替代梯度停止条件。报告表中注明未达到精度或发散的试验，不能将较早发散误判为收敛更快。新值没有优于原值，只要选择与分析合理，仍可得满分。本节效率最优倍数不直接等于最终提交倍数；第 2.2 节按验证误差联合选择最终学习率与正则化参数。不得通过测试集选学习率。


\Needspace{5\baselineskip}\subsection*{2.2 调整正则化参数（12 分）}
固定第 0.2 节的设备划分。先完成第 2.1 节，再用梯度下降联合比较特征方案、正则化参数与学习率。每个候选从零初始化，最多更新 200 次；梯度各分量绝对值的最大值不超过 $10^{-7}$ 时可提前停止，发散保护沿用第 1.4 节。用停止时的参数评价，不允许用直接解替代或修正梯度下降结果。

对当前候选的训练特征矩阵 $Z$、训练样本数 $n$ 和正则化参数 $\lambda$，脚本重新计算


\[
L=\lambda_{\max}(Z^\top Z/n)+\lambda,\qquad
\eta=\frac{c}{L}.
\]
其中 $c$ 从第 2.1 节试过的三个稳定范围内倍数中取值：$0.01$、$1$ 和你补选的 $c_{\mathrm{new}}$；$2.1$ 仅用于发散观察，不进入模型选择。更换特征、正则化参数或重训数据后都重新计算 $L$，不能沿用上一阶段的实际学习率。若 $L=0$，脚本取 $\eta=c$，在初始梯度为零时停止。

\textbf{（a）第一轮粗比较（4 分）。} 运行以下命令，对两种特征、三个学习率倍数与以下三个正则化参数的组合进行训练，共 18 个梯度下降候选：


\[
\lambda\in\{0,\ 0.1,\ 1\}.
\]

\par\addvspace{0.65\baselineskip}\noindent\begin{minipage}{\linewidth}
\begin{lstlisting}[aboveskip=0pt,belowskip=0pt]
python src/run.py --stage compare --data data --out results
\end{lstlisting}
\end{minipage}\par\addvspace{0.65\baselineskip}
脚本自动完成组合训练。\code{coarse_joint_comparison.csv} 保留 18 行完整结果；\code{coarse_comparison.csv} 对每个“特征与正则化参数”组合保留验证误差最低的学习率，再加训练标签均值对照，共 7 行。均值对照不参与最终配置选择。查看训练与验证 RMSE、系数范数和迭代状态，决定下一轮探索方向。不能根据本地测试成绩决定方向。

\textbf{（b）自行补选两个值（4 分）。} 根据第一轮验证结果，选两个互不相同的新正则化参数，不得重复已有值。报告选择方向与理由，再通过 \code{--extra-lambdas} 传入；脚本自动在两种特征和三个学习率倍数下运行，无需增加手工试验轮数。下面仅为命令格式示例：


\par\addvspace{0.65\baselineskip}\noindent\begin{minipage}{\linewidth}
\begin{lstlisting}[aboveskip=0pt,belowskip=0pt]
python src/run.py --stage compare --extra-lambdas 0.02 0.05 --out results
\end{lstlisting}
\end{minipage}\par\addvspace{0.65\baselineskip}
\textbf{（c）冻结最终配置（4 分）。} 第二轮保留原有组合，\code{joint_comparison.csv} 共 30 行；按每个特征与正则化参数组合保留最佳学习率，再加均值对照，得到 11 行 \code{comparison.csv} 和 660 行 \code{validation_predictions.csv}。报告两轮汇总表、最终三项配置，并从完整表中引用同一特征、同一正则化参数下不同学习率的验证误差，说明有限训练预算的影响。无需重复粘贴全部 30 行。

\code{regularization_curves.png} 固定 $c=1$，显示正则化参数变化时的训练与验证 RMSE、系数范数；不能把同时改变学习率的差异全部归因于正则化。预算内的参数未必达到目标最优解，因此不要求误差或范数严格单调。

脚本在 30 个梯度下降候选中按未舍入验证 RMSE 选择最终配置；精确相同时依次优先线性特征、较大的正则化参数、较小的学习率倍数。结果写入 \code{config.json}，不要手工改写。完成学习率试验后再开始本节，模型比较开始后不得返回修改学习率候选；不得看过测试成绩后重新选择。新增值不必优于第一轮结果，不要求继续无上限搜索。


\Needspace{5\baselineskip}\subsection*{2.3 两项评估对照（4 分）}
\code{compare} 命令同时生成三行 \code{audit.csv}。三种条件都固定二次特征、正则化参数 $\lambda=0.1$、学习率倍数 $c=1$ 和 200 次梯度下降更新预算，不在对照中另行调参。


{\small\renewcommand{\arraystretch}{1.18}
\begin{longtable}{@{}>{\raggedright\arraybackslash}p{\dimexpr 0.26\linewidth-2\tabcolsep\relax}>{\raggedright\arraybackslash}p{\dimexpr 0.74\linewidth-2\tabcolsep\relax}@{}}
\toprule
\textbf{条件} & \textbf{划分与特征} \\
\midrule
\endfirsthead
\toprule
\textbf{条件} & \textbf{划分与特征} \\
\midrule
\endhead
\midrule
\multicolumn{2}{r}{\footnotesize 接下页}\\
\endfoot
\bottomrule
\endlastfoot
group & 原设备划分，480 条训练、60 条验证，32 个二次特征 \\
random & 仅将上述 540 条训练及验证记录重新按行打乱为 480、60 条；特征同 group \\
post & 划分同 group，额外加入运行结束后的 \code{post_cycle_meter} \\
\end{longtable}}
三种条件均不使用已保留的 60 条本地测试记录。\code{random} 使用相同随机种子，其训练和验证中可能出现相同设备。\code{post} 是指定的错误信息使用示例，不能成为最终模型的候选。

\textbf{你需要做：}报告三个 RMSE，以及 \code{random} 相对 \code{group}、\code{post} 相对 \code{group} 的 RMSE 差；报告脚本统计的训练与验证共享设备数。回答第 4 部分对应问题。不要求任何对照一定取得更低误差，也不要求显著性检验。

计分：同模型与同参数的三行对照完整 2 分，两项误差差值正确 1 分，共享设备数正确 1 分。解释的分数另计入报告，不重复扣分。


\Needspace{5\baselineskip}\subsection*{2.4 测试与保存（4 分）}
完成两项调参及对照分析后，把所选配置写入 \code{src/submission.py} 的 \code{get_config()}。此接口已提供，只需填写三个值，不增加计算函数或分值：


\par\addvspace{0.65\baselineskip}\noindent\begin{minipage}{\linewidth}
\begin{lstlisting}[aboveskip=0pt,belowskip=0pt]
def get_config():
    return {
        "basis": None,
        "lambda": None,
        "lr_factor": None,
    }
\end{lstlisting}
\end{minipage}\par\addvspace{0.65\baselineskip}
将 \code{basis}、\code{lambda}、\code{lr_factor} 分别填写为 \code{config.json} 中的 \code{basis}、\code{lam}、\code{lr_factor}。三项均来自联合验证选择，不要填写第 2.1 节的 \code{chosen_factor} 或实际学习率。应原样复制所选数值，不做四舍五入。

\code{basis} 只允许 \code{linear} 或 \code{quadratic}；训练标签均值仅作为对照，不允许提交为最终模型。正则化参数须为有限非负数，学习率倍数须为有限数且在开区间 $(0,2)$ 内。接口必须返回上述三个字段，不得读取本地配置文件、结果目录或执行参数搜索；提交前不能保留 \code{None}。

最终训练调用 \code{train_gd()}，使用所选特征、正则化参数和学习率倍数；本地测试重训与 OJ 最终训练均从零开始、最多更新 200 次、使用相同梯度停止阈值。学习率和正则化参数共同影响预算内得到的系数及最终性能。不得增加预算、热启动、替换为直接解或伪造训练历史。直接求解函数仅用于第 1 部分核验和第 2.1 节的精度参照。

冻结代码中的配置后，运行：


\par\addvspace{0.65\baselineskip}\noindent\begin{minipage}{\linewidth}
\begin{lstlisting}[aboveskip=0pt,belowskip=0pt]
python src/run.py --stage final --data data --out results
\end{lstlisting}
\end{minipage}\par\addvspace{0.65\baselineskip}
脚本检查已经完成学习率补选与正则化参数补选，核对 \code{get_config()} 与两份选参记录一致；不一致则在测试前报错。通过后，脚本读取该接口，保存 \code{submitted_config.json}，再按第 0.2 节重训、测试、用全部数据重训，生成 \code{test_metrics.csv}、60 行 \code{test_predictions.csv} 和 \code{final_model.npz}。报告三个本地测试指标，确认 {\small\texttt{run.predict\_saved(model\_path, raw\_X)}} 可以加载模型。生成的结果和模型文件仅留存本地，不上传。计分：配置冻结与正确重训 2 分，模型可加载且预测一致 2 分。

已存在测试结果时，脚本阻止再次调参或重复测试。正常调试应在测试前完成；若发现导致测试无效的程序错误，可修正后如实记录原因及重跑过程，不得利用已见测试成绩修改参数。


\Needspace{7\baselineskip}\section{选做探索（不计分）}
可以另与 scikit-learn 的 Ridge 对照，或研究特征尺度对梯度下降的影响。选做结果单独标注，不替换规定的两轮调参，不依据本地测试结果做探索。不要求实现 Lasso、分类器、重复抽样稳定性分析或偏差与方差的独立模拟。


\Needspace{7\baselineskip}\section{报告问题（16 分）}
按题号回答以下四个问题，每题 4 分。每题建议一段文字，引用自己的结果；可以附一条必要公式，不需要写长篇理论综述。

\textbf{（1）预处理。} 为什么验证集必须使用训练集的均值和标准差？小样例中的常数列如何处理？结合你保存的检查结果说明实现是否符合这一要求。

\textbf{（2）训练。} 写出梯度公式，结合有限差分检查及四条学习率曲线，解释小步长、大步长与所选新步长的表现。为什么不能只看“损失下降”或最终迭代次数判断训练正确？区分学习率选择与正则化参数选择的作用。

\textbf{（3）正则化。} 结合两轮结果，解释你为何选择两个新参数。训练误差、验证误差和系数范数如何变化？选中配置是否改变？说明为何系数更小或训练更充分不意味着验证误差必然更低，不把一次划分的结果当成普遍规律。

\textbf{（4）评估。} 结合第 2.3 节的差值与共享设备数，解释按行划分和加入事后字段分别改变了什么。哪种条件与“运行前预测新设备”一致？说明为什么设备隔离不能修复事后字段问题，以及为什么单次对照不能精确量化所有场景下的乐观程度。


\Needspace{7\baselineskip}\section{流程记录与反馈（4 分）}
报告正文建议 4--6 页：简述流程，列出四组学习率试验、两轮参数比较、三个测试指标和三行对照；放入 \code{lr_curves.png} 与 \code{regularization_curves.png}，再回答四个问题。初始检查曲线可放附录，第一轮结果可在最终 11 行表中标注，无需重复贴表。必须保留调参前的观察与选择理由，不要求长篇推导。

流程与运行记录 2 分：写清环境、命令、配置选择和重训顺序。AI 使用或自检记录 1 分：写明使用工具及用途，选一个关键建议或假设说明如何通过运行独立核验；未使用 AI 则记录等价自检。反馈 1 分：如实填写完成实验的大致用时，可另提课程建议。

允许使用 AI 辅助，但必须理解并能解释提交的代码、图表和结果。总分 100 分。OJ 系统负责代码 40 分与最终性能 10 分的评分，共 50 分；Blackboard 实验报告占其余 50 分，包括调参与实验结果 30 分、四个报告问题 16 分、流程记录与反馈 4 分。环境创建不计分；最终性能按第 6 节与固定基准比较。实验报告以 PDF 格式上传至 Blackboard 系统（\href{https://bb.ustc.edu.cn}{bb.ustc.edu.cn}），\code{submission.py} 单独上传至 OJ（\href{https://oj.temaurinum.moe}{oj.temaurinum.moe}）；不接收其他文件。图表、参数选择和运行记录写入报告，不另交结果目录、模型文件或压缩包。报告须记录补选学习率、两个补选正则化参数、最终配置和完整运行命令，以便复现。报告须与代码中的配置一致。OJ 的代码 40 分独立测试七个计算函数，不因最终特征方案而免测。


\Needspace{7\baselineskip}\section{最终性能评测（10 分）}
本项 10 分由 OJ 系统评分。OJ 仅从 \code{submission.py} 读取 \code{get_config()}，在全部开发数据上调用提交的特征、预处理和 \code{train_gd()}，按第 2.2 节的 200 次预算与停止条件训练模型，学习率与正则化参数均来自该配置。不读取报告、配置文件或已保存模型。OJ 以 RMSE 将所得模型与固定的 scikit-learn Ridge 基准（二次特征，正则化参数 $\lambda=0.1$）比较。RMSE 不高于基准得 10 分；否则，记基准与学生模型的 RMSE 分别为 $E_b$ 和 $E_s$，性能得分 $P$ 按以下比例计算：


\[
P=10\frac{E_b}{E_s}.
\]
\code{final_model.npz} 仅在本地保存，无需上传。\code{performance_preview.csv} 仅为本地预估，不作为正式得分，也不能据此继续调参。


\end{document}
